\subsection{Présentation de PHP et Laravel}

\textbf{PHP} est un langage de script côté serveur largement déployé sur le web.
Historiquement associé au développement de sites dynamiques, il a évolué vers des
versions modernes (8.x) intégrant un typage strict, des performances accrues (JIT) et
un écosystème mature de frameworks.

\medskip

\textbf{Laravel}~\cite{laravel} est le framework PHP le plus adopté pour construire des
applications web et des API. Créé par Taylor Otwell, il impose une structure claire
autour du patron \textbf{MVC} (\textit{Model-View-Controller}) et propose des
composants intégrés pour les tâches courantes : routage, validation, authentification,
files d'attente, envoi de courriels et gestion des fichiers.

\subsection{Architecture MVC dans Laravel}

Le patron \textbf{MVC} répartit les responsabilités :

\begin{itemize}
  \item \textbf{Modèle} : représentation des entités métier et accès aux données
  (Eloquent).
  \item \textbf{Vue} : dans une API, les \textit{Resources} et \textit{Collections}
  formatent les réponses JSON pour le client.
  \item \textbf{Contrôleur} : reçoit la requête HTTP, applique la validation, invoque
  les services métier et retourne la réponse.
\end{itemize}

Cette organisation facilite la lecture du code et l'ajout de nouveaux modules. Pour
l'application CIVCO BTP, chaque domaine métier (projets, clients, devis, factures)
dispose de son contrôleur, de ses modèles et de ses règles de validation.

\subsection{Composants clés de l'écosystème Laravel}

\subsubsection{Routage et validation}

Le fichier de routes (\texttt{routes/api.php}) déclare les points d'entrée de l'API.
Chaque route est associée à un contrôleur et protégée par des \textit{middleware}
(vérification d'authentification, contrôle de permissions, contexte entreprise).

Les \textbf{Form Requests} encapsulent les règles de validation des données entrantes
(ex.~: champs obligatoires d'un devis, formats de dates, montants positifs), garantissant
l'intégrité des données avant persistance.

\subsubsection{Services et logique métier}

Au-delà des contrôleurs, la logique complexe est extraite dans des classes
\textbf{Service} (\texttt{QuoteService}, \texttt{PrintTrackingService}, etc.). Ce
découpage respecte le principe de responsabilité unique : le contrôleur orchestre,
le service implémente la règle métier.

\subsubsection{Migrations et seeders}

Les \textbf{migrations} versionnent le schéma de la base de données sous forme de fichiers
PHP exécutables de façon reproductible. Les \textbf{seeders} initialisent des données
de référence (rôles, permissions, données de démonstration), indispensables pour
tester l'application en conditions proches du réel.

\subsection{Comparaison avec les alternatives PHP}

\begin{table}[H]
\caption{Comparaison des frameworks PHP pour une API REST}
\label{tab:php_frameworks}
\centering
\small
\begin{tabular}{|l|p{4cm}|p{4cm}|p{4cm}|}
\hline
\textbf{Framework} & \textbf{Philosophie} & \textbf{Forces} & \textbf{Contexte idéal} \\
\hline
\textbf{Laravel} & Convention over configuration & Écosystème complet, Sanctum, Eloquent &
API + SPA, PME/ETI \\
\hline
Symfony & Composants réutilisables & Modularité, robustesse & Grandes applications \\
\hline
CodeIgniter & Minimaliste & Légèreté & Petits projets \\
\hline
\end{tabular}
\end{table}

Laravel ressort du Tableau~\ref{tab:php_frameworks} pour un projet de gestion métier
nécessitant authentification, RBAC, documents et évolutivité --- le profil exact du
projet CIVCO BTP.

\subsection{Laravel dans le projet CIVCO BTP}

L'API Laravel expose plus de quarante contrôleurs organisés par domaine, des
\textit{middleware} de contrôle d'accès et un ensemble de services métier. Cette
structure permet d'ajouter progressivement de nouveaux modules (avenants contractuels,
modèles de documents, suivi d'impression) sans remettre en cause l'architecture
existante.

\subsection{Middleware, Resources et couche de présentation API}

\subsubsection{Middleware}

Les \textbf{middleware} Laravel interceptent chaque requête HTTP avant qu'elle
n'atteigne le contrôleur. Dans le projet, plusieurs middlewares assurent la
sécurité et le contexte métier :

\begin{itemize}
  \item \textbf{Authentification Sanctum} : vérifie que l'utilisateur est connecté.
  \item \textbf{Contrôle de permissions} : valide que l'utilisateur possède la
  permission requise (\texttt{project.view}, \texttt{quote.manage}, etc.).
  \item \textbf{Contexte entreprise} : identifie la société courante via l'en-tête
  \texttt{X-Company-Id}, garantissant l'isolation des données métier.
\end{itemize}

\subsubsection{API Resources}

Les classes \textbf{Resource} (\texttt{ProjectResource}, \texttt{InvoiceResource},
etc.) transforment les modèles Eloquent en tableaux JSON normalisés. Elles permettent
de contrôler précisément les champs exposés à l'interface, d'inclure conditionnellement
des relations (\texttt{client}, \texttt{lines}, \texttt{teamMembers}) et d'uniformiser
le format de réponse sur toute l'API.

\subsubsection{Validation et gestion des erreurs}

Lorsqu'une requête échoue la validation, Laravel retourne une réponse \texttt{422}
avec le détail des champs en erreur. Le frontend React intercepte ces réponses pour
afficher des messages à l'utilisateur, améliorant l'expérience de saisie sur les
formulaires de devis, factures et fiches projets.
